\documentclass[a4paper,10pt]{article}

% =========================================================
% PACKAGES
% =========================================================

% Lab 3 packages
\usepackage{amsmath}
\usepackage{algorithm}
\usepackage{algpseudocode}
\usepackage[retainorgcmds]{IEEEtrantools}

% General formatting
\usepackage[margin=1.5cm]{geometry}
\usepackage{longtable}
\usepackage{booktabs}
\usepackage{hyperref}

% If your images are inside a folder called Images,
% you can uncomment the following line:
% \graphicspath{{Images/}}

\title{\textbf{TPT1201 Research Methodology}\\
LaTeX Command Cheat Sheet\\
\large Lab 3}

\author{}
\date{}

\begin{document}

\maketitle
\tableofcontents
\newpage

% =========================================================
% INTRODUCTION
% =========================================================

\section{Introduction}

This document summarizes the important LaTeX commands and examples
covered in Lab 3.

\begin{itemize}
    \item \textbf{Lab 3:} Mathematical formulae, equations, arrays,
    matrices, piecewise functions, fractions, sums, products,
    mathematical fonts, and algorithms.
\end{itemize}

\textbf{Tip:} In Overleaf, create a new project, open \texttt{main.tex},
delete the existing contents, and paste this entire document.

\newpage

% =========================================================
% LAB 3
% =========================================================

\section{Lab 3 -- Mathematics and Algorithms}

\subsection{Learning Objectives}

Lab 3 focuses on creating mathematical formulae, simple mathematical
equations, and algorithms.

% ---------------------------------------------------------
% LAB 3 CHEAT SHEET
% ---------------------------------------------------------

\subsection{Lab 3 Command Cheat Sheet}

\renewcommand{\arraystretch}{1.4}

\begin{longtable}{p{0.25\textwidth} p{0.32\textwidth} p{0.35\textwidth}}
\caption{Lab 3 -- LaTeX Command Cheat Sheet}
\\
\toprule
\textbf{Command / Code} &
\textbf{Purpose} &
\textbf{Example / Meaning}
\\
\midrule
\endfirsthead

\toprule
\textbf{Command / Code} &
\textbf{Purpose} &
\textbf{Example / Meaning}
\\
\midrule
\endhead

\bottomrule
\endfoot

\texttt{\textbackslash usepackage\{amsmath\}} &
Loads the AMS mathematics package. &
Needed for many mathematical environments and commands.
\\

\texttt{\$ ... \$} &
Creates inline mathematics. &
\texttt{\$a+b=c\$} produces $a+b=c$.
\\

\texttt{\textbackslash begin\{equation\}} &
Starts a numbered equation. &
Used when the equation should automatically receive a number.
\\

\texttt{\textbackslash end\{equation\}} &
Ends the numbered equation environment. &
Works together with \texttt{\textbackslash begin\{equation\}}.
\\

\texttt{\textbackslash begin\{equation*\}} &
Starts an unnumbered equation. &
Useful when an equation number is not required.
\\

\texttt{\textasciicircum} &
Creates a superscript. &
\texttt{a\textasciicircum2} $\rightarrow a^2$.
\\

\texttt{\_} &
Creates a subscript. &
\texttt{x\_1} $\rightarrow x_1$.
\\

\texttt{\textbackslash nonumber} &
Prevents an equation line from being numbered. &
Useful inside multi-line equation environments.
\\

\texttt{\textbackslash begin\{IEEEeqnarray\}\{rCl\}} &
Creates an aligned mathematical equation using
IEEEeqnarray. &
Useful for equations spread over multiple lines.
\\

\texttt{\&} &
Defines alignment positions. &
Used in IEEEeqnarray to align the equals sign.
\\

\texttt{\textbackslash\textbackslash} &
Starts a new line. &
Used to continue an equation on the next line.
\\

\texttt{\textbackslash :} &
Adds mathematical spacing. &
Example: \texttt{+\:l} gives additional spacing.
\\

\texttt{\textbackslash left( ... \textbackslash right)} &
Creates automatically sized parentheses. &
Useful for matrices and large expressions.
\\

\texttt{\textbackslash begin\{array\}\{ccc\}} &
Creates an array with three centered columns. &
Used to construct matrices or tabular mathematical structures.
\\

\texttt{\textbackslash vdots} &
Creates vertical dots. &
Used when showing continuation of rows.
\\

\texttt{\textbackslash ddots} &
Creates diagonal dots. &
Commonly used in matrices.
\\

\texttt{\textbackslashmathbf\{X\}} &
Creates bold mathematical symbols. &
Example: $\mathbf{X}$.
\\

\texttt{\textbackslash begin\{cases\}} &
Common way to create piecewise functions. &
Each row represents a condition.
\\

\texttt{\textbackslash text\{...\}} &
Places normal text inside mathematics. &
Example: \texttt{\textbackslash text\{if \} }.
\\

\texttt{\textbackslash frac\{a\}\{b\}} &
Creates a fraction. &
\texttt{\textbackslash frac\{1\}\{3\}} $\rightarrow \frac{1}{3}$.
\\

\texttt{\textbackslash sum\_\{i=1\}\^n} &
Creates a summation symbol with limits. &
$\displaystyle\sum_{i=1}^{n}$.
\\

\texttt{\textbackslash prod\_\{i=1\}\^n} &
Creates a product symbol with limits. &
$\displaystyle\prod_{i=1}^{n}$.
\\

\texttt{\textbackslash mathrm\{...\}} &
Uses upright/roman mathematical font. &
$\mathrm{ABCDE}$.
\\

\texttt{\textbackslash mathit\{...\}} &
Uses italic mathematical font. &
$\mathit{ABCDE}$.
\\

\texttt{\textbackslash mathnormal\{...\}} &
Uses normal mathematical font. &
$\mathnormal{ABCDE}$.
\\

\texttt{\textbackslash mathcal\{...\}} &
Uses calligraphic mathematical font. &
$\mathcal{ABC}$.
\\

\texttt{\textbackslash usepackage\{algorithm\}} &
Loads the algorithm floating environment. &
Required for the \texttt{algorithm} environment.
\\

\texttt{\textbackslash usepackage\{algpseudocode\}} &
Loads commands for writing pseudocode. &
Provides commands such as \texttt{\textbackslash If},
\texttt{\textbackslash While}, and \texttt{\textbackslash State}.
\\

\texttt{\textbackslash begin\{algorithm\}} &
Starts an algorithm. &
Creates a numbered algorithm block.
\\

\texttt{\textbackslash caption\{...\}} &
Gives an algorithm a title. &
Example: \texttt{\textbackslash caption\{MyTestingProcedure\}}.
\\

\texttt{\textbackslash begin\{algorithmic\}[1]} &
Starts the pseudocode section with line numbering. &
The \texttt{[1]} enables line numbers.
\\

\texttt{\textbackslash Procedure\{...\}\{...\}} &
Defines a procedure. &
Example: \texttt{\textbackslash Procedure\{FindPathBK\}\{v,u,p\}}.
\\

\texttt{\textbackslash If\{...\}} &
Starts an IF condition. &
Used for conditional logic.
\\

\texttt{\textbackslash Else} &
Defines the alternative branch of an IF statement. &
Executed when the IF condition is false.
\\

\texttt{\textbackslash EndIf} &
Ends an IF block. &
Marks the end of the condition.
\\

\texttt{\textbackslash State} &
Creates a pseudocode statement. &
Example: \texttt{\textbackslash State \$p\textbackslash gets G\$}.
\\

\texttt{\textbackslash gets} &
Represents assignment in pseudocode. &
$p\gets G$ means assign $G$ to $p$.
\\

\texttt{\textbackslash While\{...\}} &
Starts a WHILE loop. &
Repeats statements while the condition is true.
\\

\texttt{\textbackslash EndWhile} &
Ends a WHILE loop. &
Marks the end of the loop.
\\

\texttt{\textbackslash EndProcedure} &
Ends a procedure. &
Marks the end of the algorithm procedure.
\\

\end{longtable}


% ---------------------------------------------------------
% LAB 3 EXAMPLES
% ---------------------------------------------------------

\subsection{Lab 3 Examples}

\subsubsection{1. Simple Mathematical Expression}

Use dollar signs for inline mathematics:

\begin{verbatim}
$c$ equals $a$ plus $b$.

Mathematically:
$c = a + b$.
\end{verbatim}

Result:

$c = a+b$.

% ---------------------------------------------------------

\subsubsection{2. Numbered Equation}

\begin{verbatim}
\begin{equation}
a^2+b^2=c^2
\end{equation}
\end{verbatim}

Result:

\begin{equation}
a^2+b^2=c^2
\end{equation}

A second equation is automatically numbered:

\begin{verbatim}
\begin{equation}
x^3+y^3=z^3
\end{equation}
\end{verbatim}

\begin{equation}
x^3+y^3=z^3
\end{equation}

% ---------------------------------------------------------

\subsubsection{3. IEEEeqnarray}

The Lab 3 material uses the \texttt{IEEEeqnarray} environment
for aligned equations.

\begin{verbatim}
\begin{IEEEeqnarray}{rCl}
a & = & b+c \\
  & = & d+e+f+g+h+i+j+k \nonumber \\
  && +\:l+m+n \\
  & = & p+q+r+s
\end{IEEEeqnarray}
\end{verbatim}

Result:

\begin{IEEEeqnarray}{rCl}
a & = & b+c \\
  & = & d+e+f+g+h+i+j+k \nonumber \\
  && +\:l+m+n \\
  & = & p+q+r+s
\end{IEEEeqnarray}

\textbf{Important:} Lab 3 specifies that
\texttt{IEEEtrantools} must be included and mentions that
\texttt{IEEEtranstools.sty} may need to be placed in the same folder
as the TEX file.

% ---------------------------------------------------------

\subsubsection{4. Array / Matrix}

\begin{verbatim}
\begin{equation*}
\mathbf{X}= \left(
\begin{array}{ccc}
x_1 & x_2 & \ldots \\
x_3 & x_4 & \ldots \\
\vdots & \vdots & \ddots
\end{array}
\right)
\end{equation*}
\end{verbatim}

Result:

\begin{equation*}
\mathbf{X}= \left(
\begin{array}{ccc}
x_1 & x_2 & \ldots \\
x_3 & x_4 & \ldots \\
\vdots & \vdots & \ddots
\end{array}
\right)
\end{equation*}

\textbf{Remember:}

\begin{itemize}
    \item \texttt{c} = center alignment
    \item \texttt{l} = left alignment
    \item \texttt{r} = right alignment
    \item \texttt{\&} = separates columns
    \item \texttt{\textbackslash\textbackslash} = new row
\end{itemize}

% ---------------------------------------------------------

\subsubsection{5. Piecewise Function}

\begin{verbatim}
\begin{equation*}
x = \left\{
\begin{array}{rl}
-x & \text{if } x<0,\\
0  & \text{if } x=0,\\
x  & \text{if } x>0.
\end{array}
\right.
\end{equation*}
\end{verbatim}

Result:

\begin{equation*}
x = \left\{
\begin{array}{rl}
-x & \text{if } x<0,\\
0  & \text{if } x=0,\\
x  & \text{if } x>0.
\end{array}
\right.
\end{equation*}

% ---------------------------------------------------------

\subsubsection{6. Fractions}

\begin{verbatim}
\begin{equation*}
P = \frac{1}{3} \times \frac{x}{y^2}
\end{equation*}
\end{verbatim}

Result:

\begin{equation*}
P = \frac{1}{3} \times \frac{x}{y^2}
\end{equation*}

% ---------------------------------------------------------

\subsubsection{7. Sum and Product}

Sum:

\begin{verbatim}
\begin{equation*}
X = \sum_{i=1}^n (x^2-x)
\end{equation*}
\end{verbatim}

\begin{equation*}
X = \sum_{i=1}^n (x^2-x)
\end{equation*}

Product:

\begin{verbatim}
\begin{equation*}
Y =
\left[
\prod_{i=1}^n (X|x_n)
\right]^{\frac{1}{2}}
\end{equation*}
\end{verbatim}

\begin{equation*}
Y =
\left[
\prod_{i=1}^n (X|x_n)
\right]^{\frac{1}{2}}
\end{equation*}

% ---------------------------------------------------------

\subsubsection{8. Mathematical Fonts}

\begin{verbatim}
$\mathrm{ABCDE \qquad abcde \qquad 12345}$

$\mathit{ABCDE \qquad abcde \qquad 12345}$

$\mathnormal{ABCDE \qquad abcde \qquad 12345}$

$\mathcal{ABCDE}$
\end{verbatim}

Examples:

\[
\mathrm{ABCDE \qquad abcde \qquad 12345}
\]

\[
\mathit{ABCDE \qquad abcde \qquad 12345}
\]

\[
\mathnormal{ABCDE \qquad abcde \qquad 12345}
\]

\[
\mathcal{ABCDE}
\]

% ---------------------------------------------------------
% ALGORITHMS
% ---------------------------------------------------------

\subsection{Writing Algorithms}

Lab 3 uses:

\begin{verbatim}
\usepackage{algorithm}
\usepackage{algpseudocode}
\end{verbatim}

\subsubsection{Algorithm Example 1}

\begin{verbatim}
\begin{algorithm}
\caption{MyTestingProcedure}
\begin{algorithmic}[1]

\Procedure{MyTestingProcedure}{$G$,$p$}

\If{$p\geq G$}
    \State $p\gets G$
\Else
    \If{$p+k\leq G$}
        \State $G\gets p+k$
    \EndIf
\EndIf

\EndProcedure

\end{algorithmic}
\end{algorithm}
\end{verbatim}

Actual result:

\begin{algorithm}
\caption{MyTestingProcedure}
\begin{algorithmic}[1]

\Procedure{MyTestingProcedure}{$G$,$p$}

\If{$p\geq G$}
    \State $p\gets G$
\Else
    \If{$p+k\leq G$}
        \State $G\gets p+k$
    \EndIf
\EndIf

\EndProcedure

\end{algorithmic}
\end{algorithm}

% ---------------------------------------------------------

\subsubsection{Algorithm Example 2}

\begin{verbatim}
\begin{algorithm}
\caption{FindPathBK}
\begin{algorithmic}[1]

\Procedure{FindPathBK}{$v$, $u$, $p$}

\If{$v=u$}
    \State \textbf{Write} $v$
\Else
    \State $w\leftarrow v$
    \While{$w\not=u$}
        \State \textbf{Write} $w$
        \State $w\leftarrow p(w)$
    \EndWhile
\EndIf

\EndProcedure

\end{algorithmic}
\end{algorithm}
\end{verbatim}

Actual result:

\begin{algorithm}
\caption{FindPathBK}
\begin{algorithmic}[1]

\Procedure{FindPathBK}{$v$, $u$, $p$}

\If{$v=u$}
    \State \textbf{Write} $v$
\Else
    \State $w\leftarrow v$
    \While{$w\not=u$}
        \State \textbf{Write} $w$
        \State $w\leftarrow p(w)$
    \EndWhile
\EndIf

\EndProcedure

\end{algorithmic}
\end{algorithm}

% =========================================================

\section{Final Reminder}

\begin{itemize}
    \item \textbf{Lab 3 = Mathematics + Algorithms}
    \item Always check that required packages are included in the
    preamble.
    \item In tables, \texttt{\&} separates columns and
    \texttt{\textbackslash\textbackslash} creates a new row.
    \item \texttt{\textbackslash hline} creates a full horizontal line,
    while \texttt{\textbackslash cline\{i-j\}} creates a partial line.
    \item \texttt{\textbackslash multicolumn} merges columns.
    \item \texttt{\textbackslash multirow} merges rows.
    \item \texttt{\textbackslash frac} creates fractions.
    \item \texttt{\textasciicircum} creates superscripts and
    \texttt{\_} creates subscripts.
    \item \texttt{algorithm} and \texttt{algpseudocode} are used for
    algorithms.
    \item \texttt{graphicx} is required for inserting images.
\end{itemize}

\end{document}
